problem:  
Let \( M^n \) be a closed, simply connected Riemannian manifold with nonnegative sectional curvature that admits an isometric \( S^1 \)-action which is **fixed point homogeneous**; that is, the fixed point set \(\mathrm{Fix}(S^1, M)\) contains a component \(F\) whose normal sphere in \(M\) meets every principal orbit. Assume that the fixed component \(F\) has codimension at least \(3\).  

Does it follow that \(M\) is diffeomorphic to a rank-one symmetric space — that is, to one of \(S^n\), \(\mathbb{CP}^k\), \(\mathbb{HP}^k\), or the Cayley plane \(\mathbb{CaP}^2\)?  

Why is it a "good" problem:  
This problem asks whether the deep rigidity phenomena known in **positive sectional curvature geometry** extend into the **nonnegative curvature** realm for manifolds with **fixed point homogeneous isometric circle actions**. In the positive curvature case, Grove and Searle established a complete classification: any such manifold is diffeomorphic to a rank-one symmetric space. However, when curvature is merely nonnegative, new phenomena emerge — potential metric collapses, singular orbit spaces of Alexandrov type, and more intricate topology for the fixed sets and orbit spaces.  

The **codimension ≥ 3 condition** pushes this problem beyond known classification techniques. In lower codimension, disk-bundle decompositions together with cohomogeneity-one arguments can often fully describe the manifold, but in higher codimension, such methods break down. Investigating whether nonnegative curvature retains this strong global rigidity would require developing new analytic or topological tools — possibly involving **Alexandrov geometry of the orbit space**, analysis of **metric join structures**, or insights from **Ricci flow and collapse theory**.  

Thus the problem exemplifies a "good" problem because:  
- It is **simple to state** but demands new ideas from multiple domains — Riemannian geometry, transformation groups, and global topology.  
- It **extends a foundational rigidity result** beyond its current frontier, asking whether the geometry/topology dichotomy of positively curved fixed point homogeneous spaces remains valid under weaker curvature conditions.  
- It **invites conceptual innovation**: a solution would either establish a new geometric classification principle in the nonnegative curvature setting or uncover previously unseen examples, both of which would reshape the understanding of global symmetry and curvature interaction.